\documentclass[11pt]{amsart}

\usepackage[T1]{fontenc}
\usepackage{lmodern}
\usepackage{microtype}
\usepackage{mathtools,amssymb,amsthm}
\usepackage[hidelinks]{hyperref}

\hypersetup{
  pdftitle={Lin's hot-cold maximum-determinant correspondence at order ten}
}

\newtheorem{theorem}{Theorem}
\newtheorem{lemma}[theorem]{Lemma}
\newtheorem{proposition}[theorem]{Proposition}
\newtheorem{corollary}[theorem]{Corollary}
\theoremstyle{remark}
\newtheorem*{remark}{Remark}

\newcommand{\Fk}{\mathcal F_k}
\newcommand{\Gk}{\mathcal G_k}

\title[Lin's correspondence at order ten]
{Lin's Hot--Cold Maximum-Determinant Correspondence at Order Ten}

\date{}

\subjclass[2020]{05B20, 15A15, 05C20}
\keywords{maximal determinant, skew-symmetric sign matrix, hot matrix, cold matrix,
Seidel tournament matrix, exhaustive census}

\begin{document}

\begin{abstract}
A \emph{cold} matrix is a skew-symmetric matrix with zero diagonal and off-diagonal
entries $\pm1$; the \emph{hot} matrix attached to it is $H=C+I$. Writing $f_k(n)$ and
$g_k(n)$ for the largest determinants attained in the two families of order~$n$, Lin
asked, for $n\equiv2\pmod4$, whether $\det H=g_k(n)$ holds precisely when
$\det(H-I)=f_k(n)$. We settle the cell $n=10$ by exhaustion. Bordering reduces the
$2^{45}$ labelled order-$10$ cold matrices to the $191{,}536$ isomorphism classes of
$9$-vertex tournaments, and a census of a covering by $6880\cdot2^{8}=1{,}761{,}280$
bordered matrices gives $f_k(10)=33489=183^2$ and $g_k(10)=64000=125\cdot2^{9}$, with
\emph{exactly one} switching-and-permutation class attaining either value. That the
$6880$-tournament list carrying the covering is complete and has no repeated class is not
imported: it is certified inside the verification by an exact orbit count against the
$2^{28}$ labelled $8$-vertex tournaments. Lin's
equivalence therefore holds at $n=10$ in both directions, quantified over all
maximisers. The maximiser is exhibited, and its two determinants are read off one
characteristic polynomial by hand. Both values are already reported in the literature,
and $f_k(10)$ is proved there.
\end{abstract}

\maketitle

\section{The question}\label{sec:q}

The source is a post on Peter Cameron's blog \cite{Cameron2011}, catalogued as Cameron's
Problem~104. The definitions are Cameron's:

\begin{quote}
``We define a hot matrix to be a square matrix with entries $\pm1$, with $+1$ on the
diagonal and off-diagonal entries skew-symmetic''\ (sic) ``; a cold matrix is a
skew-symmetric matrix with $0$ on the diagonal and all off-diagonal elements $\pm1$.
There is an obvious bijection between cold and hot matrices, given by $H = C+I$.''
\end{quote}

\noindent and so is the question:

\begin{quote}
``Dennis Lin asked whether it is true that, for $n$ congruent to $2 \pmod 4$, if $H$ is a
hot matrix with maximum determinant, then $H-I$ is a cold matrix with maximum
determinant. Perhaps it goes the other way as well.''
\end{quote}

Write $\Fk(n)$ for the cold matrices of order~$n$ and $\Gk(n)=\{C+I: C\in\Fk(n)\}$ for the
hot ones, and put
\[
  f_k(n)=\max_{C\in\Fk(n)}\det C,
  \qquad
  g_k(n)=\max_{H\in\Gk(n)}\det H .
\]
Odd orders are excluded because a skew-symmetric matrix of odd order is singular, and
orders divisible by~$4$ are settled by Hadamard's bound whenever a skew-Hadamard matrix
exists; ``What happens if $n$ is not a multiple of $4$? This was Lin's question''
\cite{Cameron2011}. The same post records the evidence available for the first such
order, and warns that the bijection $C\mapsto C+I$ ``does not induce a monotone map on
their determinants'':

\begin{quote}
``There are too many matrices for $n = 10$ for an exhaustive search; and in lieu of doing
anything clever, I simply looked at a few million random matrices \ldots\ (The best I
found were 64000 for hot and 33489 for cold.)''
\end{quote}

Armario and Frau restate the question as an explicit equivalence in a refereed journal
\cite[\S1]{ArmarioFrau2016}: ``it asks to decide whether or not the following equivalence
holds: $M$ has maximum determinant among the matrices in $F_k(n)$ if and only if $M + I$
has maximum determinant among the matrices in $G_k(n)$''. That is the statement we prove
at $n=10$. Their Theorem~2.6 proves the equivalence when
$f_k(n)=(2n-3)(n-3)^{(n-2)/2}$, which by their Proposition~2.4 requires $2n-3$ to be a
perfect square, so at $n=10$, where $2n-3=17$, it does not apply.

\section{The result}\label{sec:res}

\begin{theorem}\label{thm:main}
$f_k(10)=33489=183^2$ and $g_k(10)=64000=125\cdot2^{9}$. Moreover exactly one
switching-and-permutation class of order-$10$ cold matrices attains $f_k(10)$, exactly one
attains $g_k(10)$, and the two classes are the same one, namely that of $C=H-I$ for the
matrix $H$ displayed in \eqref{eq:H} below. Consequently, for $n=10$, a hot matrix has
maximum determinant if and only if the cold matrix $H-I$ has maximum determinant.
\end{theorem}

The witness, given as the hot matrix $H=C+I$: rows run top to bottom, \texttt{+} is $+1$
and \texttt{-} is $-1$; $H$ is hot, i.e.\ its diagonal is $+1$ and $h_{ij}=-h_{ji}$ off
the diagonal, and $C=H-I$ is cold.
\begin{equation}\label{eq:H}
\begin{array}{l}
\texttt{++++++++++}\\
\texttt{-+++++--+-}\\
\texttt{--+++++--+}\\
\texttt{---++-++--}\\
\texttt{----++-+++}\\
\texttt{---+-++-+-}\\
\texttt{-+--+-+-++}\\
\texttt{-++--+++--}\\
\texttt{--++---+++}\\
\texttt{-+-+-+-+-+}
\end{array}
\end{equation}

That this matrix attains both values is a hand computation. The characteristic polynomial
of $C$ is
\begin{equation}\label{eq:cp}
  \det(xI-C)=x^{10}+45x^{8}+770x^{6}+6210x^{4}+23485x^{2}+33489 .
\end{equation}
Its constant term is $\det C=33489=183^2$, and since the degree is even,
$\det H=\det(I+C)$ is the value of \eqref{eq:cp} at $x=-1$, that is
\[
  1+45+770+6210+23485+33489=64000=125\cdot2^{9}.
\]
The coefficient $45$ is forced and needs no computation: $\operatorname{tr}(C^{\mathsf T}C)=90$
because all $90$ off-diagonal entries have modulus~$1$, and that trace is twice the
$x^{8}$~coefficient. The sorted out-degree sequence of the tournament $\{i\to j:
h_{ij}=+1\}$ is $3,3,4,4,4,4,4,5,5,9$; as an undirected graph on $\{0,\dots,9\}$ with an
edge $\{i,j\}$, $i<j$, exactly when $h_{ij}=+1$, the same object is the $31$-edge graph
with \texttt{graph6} string \verb|I~~lkrMig|.

\section{The proof}\label{sec:proof}

Only the maximality and the uniqueness remain, and they come from four lemmas and a
finite census.

\begin{lemma}[invariance]\label{lem:inv}
Let $C$ be cold of order~$n$, let $D$ be diagonal with entries $\pm1$ and let $P$ be a
permutation matrix. Then $DCD$ and $PCP^{\mathsf T}$ are cold, and all three of
$\det C$, $\det(I+C)$ and the polynomial $\det(xI-C)$ are unchanged.
\end{lemma}

\begin{proof}
Both maps preserve the zero diagonal and the moduli of the entries, and skewness is
preserved because $D$ and $P$ are real. Since $D=D^{-1}$ and $P^{\mathsf T}=P^{-1}$, both
maps are similarities, so they fix the characteristic polynomial; and
$I+DCD=D(I+C)D$, $I+PCP^{\mathsf T}=P(I+C)P^{\mathsf T}$, so they fix $\det(I+C)$ as well.
\end{proof}

Call $C,C'$ \emph{equivalent} when $C'=PDCDP^{\mathsf T}$ for some such $D,P$; by
Lemma~\ref{lem:inv} both determinants are invariants of the equivalence class. For a
tournament $T$ on $k$ vertices let $A(T)$ be its Seidel matrix, $a_{ij}=+1$ if $i\to j$
and $-1$ if $j\to i$, and put
\[
  \operatorname{bd}(T)=
  \begin{pmatrix} 0 & \mathbf 1^{\mathsf T}\\ -\mathbf 1 & A(T)\end{pmatrix},
\]
a cold matrix of order $k+1$.

\begin{lemma}[bordering]\label{lem:border}
Every cold matrix of order $n$ is equivalent to $\operatorname{bd}(T)$ for some tournament
$T$ on $n-1$ vertices. If $T\cong T'$ then $\operatorname{bd}(T)$ and
$\operatorname{bd}(T')$ are equivalent.
\end{lemma}

\begin{proof}
Index rows by $0,1,\dots,n-1$ and switch by $D=\operatorname{diag}(d_0,\dots,d_{n-1})$ with
$d_0=+1$ and $d_j=c_{0j}$ for $j\ge1$. The $(0,j)$ entry of $DCD$ is
$d_0c_{0j}d_j=c_{0j}^2=+1$, so row~$0$ of $DCD$ is $(0,+1,\dots,+1)$ and column~$0$ is its
negative; the remaining block is the Seidel matrix of a tournament on $n-1$ vertices. The
second assertion is the case of Lemma~\ref{lem:inv} in which $P$ fixes index~$0$: those
$P$ act on the lower block as the full symmetric group on $n-1$ letters.
\end{proof}

This is the standard bordering of Reid and Brown, used in the same form by Klanderman,
Montee, Piotrowski, Rice and Shader \cite{Klanderman2025}.

\begin{lemma}[covering]\label{lem:cover}
Let $\mathcal S$ be a set of tournaments on $n-1$ vertices meeting every isomorphism
class. Then $\{\operatorname{bd}(T):T\in\mathcal S\}$ meets every equivalence class of cold
matrices of order~$n$. In particular $\max_{T\in\mathcal S}\det\operatorname{bd}(T)=f_k(n)$
and $\max_{T\in\mathcal S}\det(I+\operatorname{bd}(T))=g_k(n)$, and if all the maximisers in
$\mathcal S$ are isomorphic tournaments then exactly one equivalence class of cold
matrices attains the corresponding maximum.
\end{lemma}

\begin{proof}
Immediate from Lemma~\ref{lem:border}, since each class contains some
$\operatorname{bd}(T)$ and $\mathcal S$ contains a tournament isomorphic to that~$T$. For
the last clause: every class attaining the maximum contains some $\operatorname{bd}(T)$
with $T\in\mathcal S$ maximal, and isomorphic $T$ give equivalent matrices by
Lemma~\ref{lem:border}, so all such classes coincide.
\end{proof}

This is what makes Cameron's ``too many matrices'' finite. At $n=10$ there are
$2^{45}=35{,}184{,}372{,}088{,}832$ labelled cold matrices, but at most
$A000568(9)=191{,}536$ equivalence classes.

The covering used below is built from a list of $8$-vertex tournaments, so Lemma~\ref{lem:cover}
is only as good as the guarantee that the list meets every isomorphism class. That guarantee is
not taken from the literature: it is certified by the program, and the certificate is the
following one.

\begin{lemma}[the class list]\label{lem:classlist}
For $k\le 8$ let $R_k$ be the list obtained from $R_1=\{K_1\}$ by forming all $2^{k-1}$
one-vertex extensions of every member of $R_{k-1}$ and keeping one tournament per value of the
canonical code
\[
  \operatorname{cd}(T)=\max_{\sigma}\,\bigl(\text{upper-triangle code of }\sigma T\bigr),
\]
where $\sigma$ runs over the relabellings that order the out-degrees of $T$
non-decreasingly. Then $|R_8|=6880$, the members of $R_8$ are
pairwise non-isomorphic, and every $8$-vertex tournament is isomorphic to exactly one of them.
\end{lemma}

\begin{proof}
Two facts, each verified computationally as stated.

\emph{(i) Nothing is lost.} $\operatorname{cd}(T)$ is by construction the upper-triangle code of
an \emph{explicit} relabelling of $T$, so two tournaments with equal code are both isomorphic to
the tournament decoded from that code, hence to each other; the program checks for each of the
$6880$ members that the maximising relabelling reproduces the code and that decoding it returns
that relabelling. Therefore the deduplication discards a tournament only in favour of an
isomorphic one. Since deleting a vertex from a $k$-vertex tournament leaves a $(k-1)$-vertex one,
and an isomorphism of the smaller tournaments extends to the larger, every class at level $k$
occurs among the one-vertex extensions of $R_{k-1}$; by induction $R_8$ meets every class.

\emph{(ii) Nothing is repeated.} A relabelling permutes out-degrees, so every automorphism fixes
each vertex's out-degree and $|\!\operatorname{Aut}(T)|$ may be counted over the out-degree-preserving
permutations alone; the program checks that this count agrees with the count over all $k!$
relabellings for all $532$ members of $R_1,\dots,R_7$. With $|\!\operatorname{Aut}|$ in hand,
\[
  \sum_{T\in R_k}\frac{k!}{|\!\operatorname{Aut}(T)|}=2^{k(k-1)/2}
  \qquad (2\le k\le 8),
\]
verified level by level, the right-hand side being the number of \emph{all} labelled $k$-vertex
tournaments, $268{,}435{,}456$ at $k=8$. By (i) the orbits of the members of $R_k$ already cover
those $2^{k(k-1)/2}$ tournaments, so a repeated class would make the left-hand side exceed the
right; equality forces the members to be pairwise non-isomorphic. In particular
$|R_8|=6880=A000568(8)$ is the number of classes, not merely an upper bound for it.
\end{proof}

\noindent Independently, the canonical code is verified to be a complete isomorphism invariant
by exhaustion at two orders: over all $1024$ labelled $5$-vertex tournaments each code class is
checked to be exactly one full $120$-relabelling orbit, and over all $32{,}768$ labelled
$6$-vertex tournaments the codes are checked to take exactly $56=A000568(6)$ values with each
class of size $720/|\!\operatorname{Aut}|$.

\begin{proposition}[the census]\label{prop:census}
Take $\mathcal S$ to consist of the $2^{8}=256$ one-vertex extensions of each of the
$6880$ members of the list $R_8$ of Lemma~\ref{lem:classlist}, so
$|\mathcal S|=1{,}761{,}280$; deleting a vertex from a $9$-vertex tournament leaves an
$8$-vertex one, which by Lemma~\ref{lem:classlist} is isomorphic to a member of $R_8$, and
that isomorphism extends, so $\mathcal S$ meets every isomorphism class of $9$-vertex
tournaments.
Over $\mathcal S$,
\[
  \max_T\det\operatorname{bd}(T)=33489,
  \qquad
  \max_T\det\bigl(I+\operatorname{bd}(T)\bigr)=64000,
\]
the two maxima are attained by the same nine members of $\mathcal S$, and those nine are
pairwise isomorphic tournaments, forming the single class of \eqref{eq:H}.
\end{proposition}

Theorem~\ref{thm:main} follows from Proposition~\ref{prop:census} by
Lemma~\ref{lem:cover}, the last clause of that lemma giving both the uniqueness and the
coincidence of the two argmax sets, and hence Lin's equivalence at $n=10$ in both
directions. The census is $1{,}761{,}280$ pairs of exact integer $10\times10$
determinants, computed in exact integer arithmetic; it is carried out by \texttt{verify.py}
in the accompanying material, and \texttt{verify.output.txt} is a recorded run of that
program.

\medskip
\noindent\textbf{What else was checked.} The same program reproduces the published cells
$(f_k,g_k)=(9,16)$, $(81,160)$ and $(2401,4096)$ at $n=4,6,8$, the first two also by brute
force over all $2^{6}$ and all $2^{15}$ labelled cold matrices; the isomorphism-class
counts $1,1,2,4,12,56,456,6880$ on $1$ to $8$ vertices against $A000568$ \cite{OEIS568};
the set of order-$10$ Seidel tournament determinants published in \cite{Klanderman2025};
and the two order-reversing order-$10$ pairs recorded by Orrick \cite{OrrickComment},
$(\det H,\det C)=(45056,21609)$ and $(46080,20449)$, which occur among the leaves.

\section{Relation to the literature, and scope}\label{sec:status}

Both records are already in print. Alvarez, Armario, Frau and Gudiel write, in the first
section of \cite{Alvarez2015}: ``For $n\leq30$, orders $18$, $22$, $30$ are unresolved.
The case $n=10$ was solved by Cameron, $g_k(10)=64000$ and $f_k(10)=33489$, using a random
search and again Lin's correspondence was confirmed.'' Moreover $f_k(10)=33489$ is
\emph{proved} in print: the complete set of determinants of order-$10$ Seidel tournament
matrices is computed in \cite{Klanderman2025}, and $33489$ is its largest element. The
census above establishes both values, the equivalence in both directions, and the
uniqueness of the maximising class, independently and by exhaustion; it establishes
nothing about priority, and we claim none. Of the two further items bearing on a possible
earlier tabulation, the first has now been read and does not overlap: Zeng and You
\cite{ZengYou2026} study, for odd $k$, the hereditary class $\mathcal D_k$ of tournaments
all of whose subtournaments have determinant at most $k^2$, characterising $\mathcal D_5$
and constructing, for each odd $k$, a tournament of determinant $k^2$ in
$\mathcal D_k\setminus\mathcal D_{k-2}$; no order is singled out, no maximum over an order
is computed, and the hot matrix $C+I$ does not occur there. The second, Seberry and
Yamada, \emph{Hadamard Matrices} (Wiley, 2020), ch.~12, \emph{More on Maximal Determinant
Matrices}, pp.~245--269, we were unable to consult. It is a survey chapter on maximal
determinants, so it may well record $f_k(10)$ or $g_k(10)$; as those two constants are in
any case credited above to \cite{Alvarez2015} and \cite{Klanderman2025}, the content
claimed here is the exhaustive proof of Lin's equivalence at $n=10$ quantified over all
maximisers, together with the uniqueness of the maximising class, and not either constant.

The scope is the cell $n=10$ only, against a question quantified over all
$n\equiv2\pmod4$. Orders $18$, $22$ and $30$ are unresolved before and after this note,
exactly as \cite{Alvarez2015} records, and the method here does not reach $n=18$: it would
require the isomorphism classes of $17$-vertex tournaments. Nothing here bears on the
unrestricted maximal-determinant problem.

\begin{thebibliography}{9}

\bibitem{Alvarez2015}
V.~Alvarez, J.~A. Armario, M.~D. Frau, and F.~Gudiel,
\emph{Determinants of $(-1,1)$-matrices of the skew-symmetric type: a cocyclic approach},
Open Math. \textbf{13} (2015), no.~1, 16--25.
\href{https://doi.org/10.1515/math-2015-0003}{doi:10.1515/math-2015-0003};
\href{https://arxiv.org/abs/1311.7250}{arXiv:1311.7250}.

\bibitem{ArmarioFrau2016}
J.~A. Armario and M.~D. Frau,
\emph{On skew E-W matrices},
J. Combin. Des. \textbf{24} (2016), no.~5, 461--472.
\href{https://doi.org/10.1002/jcd.21519}{doi:10.1002/jcd.21519}.

\bibitem{Cameron2011}
P.~J. Cameron,
\emph{A matrix problem},
Peter Cameron's Blog, 19 August 2011.
\href{https://cameroncounts.wordpress.com/2011/08/19/a-matrix-problem/}%
{cameroncounts.wordpress.com/2011/08/19/a-matrix-problem/};
catalogued as Cameron's Problem~104.

\bibitem{Klanderman2025}
J.~Klanderman, M.~Montee, A.~Piotrowski, A.~Rice, and B.~Shader,
\emph{Determinants of Seidel tournament matrices},
Linear Algebra Appl. \textbf{707} (2025), 126--151.
\href{https://doi.org/10.1016/j.laa.2024.11.011}{doi:10.1016/j.laa.2024.11.011};
\href{https://arxiv.org/abs/2406.09697}{arXiv:2406.09697}.

\bibitem{OrrickComment}
W.~Orrick,
comment of 25 August 2011 on \cite{Cameron2011},
recording two order-reversing order-$10$ pairs, and both records with ``appears to be''.

\bibitem{ZengYou2026}
J.~Zeng and L.~You,
\emph{On determinants of tournaments and the characterization of $\mathcal D_5$},
Discrete Math. \textbf{349} (2026), no.~2, 114766;
\href{https://arxiv.org/abs/2408.06992}{arXiv:2408.06992} (\emph{On determinants of
tournaments and $\mathcal D_k$}).

\bibitem{OEIS568}
The OEIS Foundation,
\emph{Sequence A000568: number of tournaments on $n$ nodes},
The On-Line Encyclopedia of Integer Sequences.
\href{https://oeis.org/A000568}{oeis.org/A000568}.

\end{thebibliography}

\end{document}
